\section{Conditional Current Exchange and the Value of a Material Step}
\subsection{Changing a current model changes two material quantities}
A receiver-based space provides a reproducible starting approximation. At a fixed current budget, we can replace one of its directions while retaining the others. This asks a local design question: does the replacement improve the material step generated by the model? The test keeps the material state, illumination, physical tangent, residual, prior term and damping fixed.

Let the old and new current spaces be $[V,A]$ and $[V,C]$. When the common-core and Schur systems are invertible, their Jacobian difference is
\begin{equation}\label{eq:exchange-difference}
 D=Y_C\Gamma_C^{-1}N_C-Y_A\Gamma_A^{-1}N_A.
\end{equation}
These are two changes relative to the same retained core. If that core is singular, stable endpoints may still exist; then we compute the endpoints directly rather than using this elimination.

In the common real material coordinates, write $\mathsf I_U=J_U^TJ_U$ for the data curvature. A current change alters both data curvature and the constant gradient term:
\begin{equation}\label{eq:information-offset}
\begin{aligned}
 \Delta\mathsf I&=J_o^TD+D^TJ_o+D^TD,\\
 h_U&=J_U^Tr+\ell,\qquad \Delta h=D^Tr.
\end{aligned}
\end{equation}
Here the normal equation is $H_Us_U=-h_U$. The data curvature describes which material changes produce a measured response. The constant term describes how those responses align with the present residual. Both determine the actual endpoint step difference:
\begin{equation}\label{eq:exchange-step}
 d=s_n-s_o=-H_n^{-1}(\Delta h+\Delta\mathsf I s_o).
\end{equation}
The formula follows by subtracting the two normal equations. It applies to the fixed unconstrained quadratic; changing an active constraint set requires its own optimality equations.

Information magnitude alone cannot select this step. Even the scalar models $J_o=1$ and $J_n=-1$ have identical data curvature but opposite steps for $r=1$, $\ell=0$ and $\Lambda=1$. More generally, $\Delta\mathsf I$ can be indefinite: changing the approximation to existing data includes interference terms, unlike adding independent measurements.

\subsection{Evaluate the displacement actually produced}
For the unchanged full quadratic, a proposed displacement $d$ from $x$ has signed gain. In current exchange, $x=s_o$ is the old endpoint material step:
\begin{equation}\label{eq:exchange-gain}
 Q=\Phi_F(x)-\Phi_F(x+d)=-g_F(x)^Td-\tfrac12d^TH_Fd.
\end{equation}
Positive gain means a better approximation to this frozen full-model material update. It does not by itself determine material truth error. The latter also contains a cross term with the unknown material error, and nonlinear reconstruction additionally requires feasible updates and step acceptance.

The scalar gain can be checked without a new full adjoint. Cache $u=r+J_Fx$ and compute $z=J_Fd$. Then
\begin{equation}\label{eq:directional-exchange-gain}
 Q=-u^Tz-(\ell+\Lambda x)^Td-\tfrac12(\|z\|^2+d^T\Lambda d).
\end{equation}
For complex data the first product is $\operatorname{Re}\langle u,z\rangle$. Each full directional action still pays for every illumination. After an accepted move, the endpoint and conditional candidate scores change; the cached output updates as $J_Fx\leftarrow J_Fx+z$.

If validated output-error bounds $e_u,e_z$ are available, the computed gain has error at most
\begin{equation}\label{eq:exchange-gain-bound}
 \epsilon_Q=e_u\|\widehat z\|+(\|\widehat u\|+\|\widehat z\|)e_z+e_ue_z+\tfrac12 e_z^2.
\end{equation}
A positive lower bound can justify acceptance. A local stopping certificate instead requires an upper bound below tolerance for every feasible move in the declared neighborhood. A shortlist with no accepted move does not establish that condition. The computations below use double-precision full directional checks; a rigorous output bound is unavailable, so their numerical acceptance and stopping status are reported separately.

\subsection{Information diagnostics and their scope}
We normalize data curvature with a fixed positive material scale $P_0$, kept separate from numerical Levenberg--Marquardt damping. For eigenvalues $\nu_i$ of $P_0^{-1/2}\mathsf I_UP_0^{-1/2}$, we record
\begin{equation}
 d_{\rm eff}=
 \sum_i\frac{\nu_i}{1+\nu_i},\qquad
 \mathcal D=\tfrac12\sum_i\log(1+\nu_i).
\end{equation}
These quantities summarize the response claimed by a model. We also measure its curvature discrepancy from the full reference. Greater $\mathcal D$ need not mean greater fidelity, and neither scalar retains the residual-dependent term in~\eqref{eq:information-offset}. The present data normalization is a prescribed numerical scale, not an independently calibrated noise covariance; no mutual-information interpretation is needed.

Small Gaussian material spaces permit complete reference diagnostics. For the larger voxel spaces, all endpoints use the same sixteen material probe directions; those information measures describe that projection only. Information diagnostics, full-reference steps and material truth remain evaluation data, separate from the reduced-reference scores used to propose exchanges.
